\chapter{Fiziksel Guvenlik}

\section*{Giris}
Fiziksel guvenlik, dijital guvenligin ilk ve vazgecilmez katmanidir. Veri merkezi, ofis alani, cihaz depolari ve kritik tesisler korunmadan siber kontroller tek basina yeterli olmaz.

\section{Fiziksel Cevre ve Tesis Guvenligi}
\subsection{Cevre Guvenligi: Sensorler, Kameralar ve Bariyerler}
Cevre guvenligi katmanli sekilde tasarlanmalidir: dis cevre bariyerleri, kamera kapsama planlari, hareket/sarsinti sensorleri ve alarm korelasyonu birlikte calismalidir.
\begin{itemize}
    \item Kritik alanlar icin 7/24 kayit ve saklama politikasi
    \item Kör nokta analizi ile kamera yerlesim optimizasyonu
    \item Alarm yorgunlugunu azaltan olay onceliklendirme
\end{itemize}

\subsection{Katmanli Fiziksel Erisim Kontrolleri ve Biyometri}
Kartli gecis, PIN ve biyometrik dogrulama birlikte kullanilarak yetkisiz giris riski azaltilir. Kritik odalar icin iki kisilik kural ve ziyaretci eskort modeli uygulanmalidir.

\section{Veri Merkezi (Data Center) Guvenligi}
\subsection{Iklimlendirme, Enerji Altyapisi ve Yangin Sondurme Sistemleri}
Veri merkezi surekliligi icin N+1 veya 2N enerji mimarisi, sicaklik/nem izleme, temiz gazli yangin sondurme ve duzenli tatbikatlar zorunludur.
\begin{table}[htbp]
\centering
\caption{Veri merkezi dayaniklilik kontrol listesi}
\begin{tabular}{|p{4cm}|p{8cm}|}
\hline
\textbf{Bilesen} & \textbf{Asgari Gereksinim} \\
\hline
Enerji & UPS + jeneratör + otomatik gecis testi \\
\hline
Iklimlendirme & Sicaklik/nem esik alarmi ve yedek klima \\
\hline
Yangin & Erken tespit + temiz gazli sondurme \\
\hline
\end{tabular}
\end{table}

\section{Donanim ve Cihaz Imhasi}
\subsection{Guvenli Veri Yok Etme (Degaussing, Shredding) ve Yasam Dongusu}
Depolama medyasinin yasam dongusu boyunca envanteri tutulmali, devir/onarim/asama sonu imha surecleri kayit altina alinmalidir.
\begin{itemize}
    \item Mantiksal silme: dogrulanabilir overwrite
    \item Degaussing: manyetik medya icin etkili yok etme
    \item Shredding: fiziksel parcala ve sertifikali imha
\end{itemize}

\Uyari{Imha sureci delil zinciri ve mevzuat uyumlulugu (KVKK vb.) acisindan belgelenmeden tamamlanmis sayilmamalidir.}
